% Root proof; D independent complete literal audit passed.
\subsection{Bounded act cones through eleven monoid elements}

\begin{lemma}\label{C:eleven-bounded-all-products}
Let $M$ be a finite monoid with $|M|\leq11$ and let $S$ be
any nonempty finite right $M$-set. If its ordinary act cone
has bounded rational homology, then every product of two
rational degree-two classes is zero. Neither a first-homology
condition nor a bound on the second-cohomology rank is required.
\end{lemma}
\begin{proof}
If $H^2(BM;\mathbb Q)=0$, every degree-two class is relative
to the monoid vertex, and the suspension quotient kills products.
Otherwise the minimal Rees ideal has at least two rows and
two columns. At most two rows give the all-order conclusion
of Theorem~\ref{C:two-row-act-cups}. We may assume at least
three rows and at least two columns. The identity lies outside
the ideal, and a nontrivial maximal group would require at
least twelve ideal elements. Thus the ideal is a rectangular
band. The only possible shapes and outside-element bounds are
\[
\begin{array}{c|c}
 \text{minimal band}&q=|M\setminus K|\\\hline
 3\times2&q\leq5\\
 4\times2&q\leq3\\
 5\times2&q\leq1\\
 3\times3&q\leq2.
\end{array}
\]
The first row is Proposition~\ref{C:three-row-five-outside-all-acts}.
The next two follow from
Lemma~\ref{C:three-outside-all-two-column-act-cups}.
For the last row, Lemma~\ref{C:two-outside-act-obstruction}
and the act-cone Tor formula make fourth rational homology
zero under the stated boundedness hypothesis. This excludes
all degree-two products in the last case as well.
\end{proof}

\begin{corollary}\label{C:small-center-projective-leaves-starlimit}
Let $C$ have outward star support with any positive number
of leaves. Suppose its central monoid $M$ has at most eleven
elements and the trivial right rational module at every leaf
monoid is projective. If $BC$ has bounded rational homology,
then every rational degree-two product on $BC$ is zero.
The inward-star assertion is dual, with left modules at the leaves.
There is no bound on the numbers of cross arrows.
\end{corollary}
\begin{proof}
Apply Proposition~\ref{C:finite-target-extreme-corner} at each
leaf. Its actual orbit quotient gives a rational cohomology-ring
equivalence to a star with the same central monoid and trivial
leaves, with nonempty finite right $M$-acts $S_i$ as its legs.
Now identify those leaf vertices. On nerves this gives exactly
the ordinary act cone on $S=\coprod_i S_i$: no nonidentity
arrow exits a leaf. On realizations, identifying these finitely
many points in a connected CW space adds only a wedge of
circles up to homotopy. This preserves bounded rational
homology, cohomology in degrees at least two, and all its
products. The preceding lemma applies to that act cone.
Pulling back the vanishing through these actual quotient maps
proves the assertion for the original star. Opposition gives
the inward version.
\end{proof}
